\documentclass[11pt]{article}
\usepackage[margin=1in]{geometry}
\usepackage{amsmath,amssymb,amsthm,mathtools}
\usepackage{enumitem}
\usepackage[hidelinks]{hyperref}

\newtheorem{theorem}{Theorem}[section]
\newtheorem{lemma}[theorem]{Lemma}
\newtheorem{proposition}[theorem]{Proposition}
\newtheorem{corollary}[theorem]{Corollary}
\theoremstyle{definition}
\newtheorem{definition}[theorem]{Definition}
\newtheorem{record}[theorem]{Record}
\newtheorem{gate}[theorem]{Gate}
\theoremstyle{remark}
\newtheorem{remark}[theorem]{Remark}
\newtheorem{firewall}[theorem]{Firewall}

\newcommand{\cA}{\mathcal A}
\newcommand{\cE}{\mathcal E}
\newcommand{\cR}{\mathcal R}
\newcommand{\cW}{\mathcal W}
\newcommand{\Pm}{P_{\mathrm{marg}}}
\newcommand{\TV}{\mathrm{TV}}
\newcommand{\NP}{\mathrm{NP}}
\newcommand{\dd}{\,\mathrm d}
\newcommand{\currentversion}{5}

\title{Renewal Yang--Mills Parallel Research Notes\\
Programme I: Nonperturbative Margin Protection\\
Programme II: Sector Tightness and Topological Transport}
\author{Parallel Yang--Mills stream}
\date{Version V\currentversion}

\begin{document}
\maketitle

\begin{abstract}
This is the single cumulative file for complementary Yang--Mills programmes
that do not duplicate the main Brillouin/SCIC/ADMC stream.  Programme I bounds,
in the source-differentiable analytic norms of the main notes, the difference
between the exact marginal RG response and the certified packet-first response.
It reaches its success stop in Section~5: the exhaustive nonlinear correction
is higher order, summable, and smaller than one half of the supplied symbolic
margin on an explicit weak-coupling domain.

Programme II is appended after that completed result.  It separates finite
centre flux, compact toron/holonomy coordinates, and integer global charge;
tests the proposed cutoff- and volume-uniform sector tail against the actual
block ledger; and stops at the first stated success or failure criterion.  The
file remains Version V5 at the user's request: new work is appended in place,
the former programme-specific filename is replaced by this umbrella name, and
no additional active version is created.  Companion research is audited at
the programme checkpoint before the final decision.
\end{abstract}

\tableofcontents

\section{V1: provenance, scope, and the exact remainder ledger}

\subsection{Sources and nonduplication boundary}

The manuscript
\texttt{Renewal\_Grand\_Tensor\_Monograph\_V074(1).tex} is used as a claim
boundary and as context for the finite colour-law and Yang--Mills interfaces.
The mathematical source for the present RG estimates is the latest cumulative
Yang--Mills note available at the start of the programme,
\texttt{Renewal\_Yang\_Mills\_Mass\_Gap\_Research\_Notes\_V74.tex}, together
with its cited archive
\texttt{Renewal\_Yang\_Mills\_Mass\_Gap\_Research\_Notes\_V65\_f8.tex}.
The supplied programme memorandum fixes the five gates and the stopping rule;
it is not treated as evidence for a theorem.

The following results are imported and will not be reproved here.
\begin{enumerate}[label=\textup{(I\arabic*)},leftmargin=4em]
\item Exact Morse--Gaussian gluing after extraction of the constant, linear,
and quadratic jets; its interior Taylor defect is
$O(\beta^{-6/5})$, while its chart-exit and Wilson-exterior terms are
stretched-exponentially small (compact V6, (6.12)--(6.19)).
\item The normalized regular-kernel defect obeys
$\|\Delta_y\|_{\rm var}\le q_{\rm ker}(\beta)$, with
$q_{\rm ker}(\beta)=O(\beta^{-6/5})$, and the collar-bad primitive obeys
$q_{\rm sat}(\beta)=O(\beta^{d_*/(2M_{\rm sat})}e^{-c\beta})$
(compact V7, (7.13)--(7.25)).
\item The triangular all-colour schedule gives every regular or bad block one
historical owner and proves an absolutely convergent connected expansion,
uniformly in volume, with four source derivatives permitted after summation
(compact V8, Theorem (8.28)--(8.32)).
\item The measured redundant coordinate has exactly zero physical marginal
leakage, the zero-theta sector is invariant, and topology-changing exits are
already bad-owned (compact V9, V11, and V12).
\item One marked connected expansion gives a uniformly bounded response of the
inverse-coupling increment, while reciprocal inversion supplies two powers of
$u=\beta^{-1}$ (compact V13, (13.3)--(13.15)).
\item On the cofinal domain the aligned stable map is a strict contraction
with some $q_{\rm st}<1$ (archived Theorem 36.14), and its nonautonomous
tracking estimate is (36.17).
\end{enumerate}

Thus this stream does not repeat the one-loop Brillouin calculation, the Morse
lemma, the forest combinatorics, the bad-field Peierls estimate, or the
packet-first response compiler.  It starts at their common output interface.

\subsection{The owner alphabet}

Fix a finite regulator, a retained zero-theta sector, and one exact triangular
block elimination.  Let $\mathcal O$ be the absolutely convergent owner algebra
of compact V8 on its smaller source polydisc.  Each primitive carries its
historical support and exactly one birth tag in
\[
 \mathfrak T=\{\mathrm{loc},\mathrm{bad},\mathrm{reset},
                 \mathrm{boundary}\}.                       \tag{1.1}
\]
Here $\mathrm{loc}$ means a regular two-jet-subtracted local primitive;
$\mathrm{bad}$ means a chart exit, Wilson exterior, collar-irregular birth, or
sector transition; $\mathrm{reset}$ means a stable complement created by the
declared reference/calibration reset; and $\mathrm{boundary}$ means a genuine
global endpoint response not represented by a local saturated owner.  The
Gaussian principal shell and every extracted finite jet are carried in the
calibrated ledger, not in $\mathcal O$.

Let a connected word $W$ be a finite word in these primitives after the V8
historical saturation and canonical rooting.  Assign $W$ to the first
applicable class in the ordered list
\[
 \mathrm{boundary}\succ\mathrm{bad}\succ\mathrm{reset}
 \succ\mathrm{polymer}\succ\mathrm{vertex},                  \tag{1.2}
\]
where ``polymer'' means that all letters are regular local letters and
$|W|\ge2$, while ``vertex'' means one regular local letter.  The order in
(1.2) is a classification rule, not a second owner: the historical owner of a
primitive is never altered.

\begin{lemma}[Disjoint exhaustive word partition]
\label{lem:word-partition}
Every non-Gaussian connected word belongs to exactly one of the five classes
in (1.2).
\end{lemma}

\begin{proof}
If a word contains a boundary tag it enters the first class.  Otherwise, if it
contains a bad tag it enters the second class; otherwise, if it contains a
reset tag it enters the third.  A remaining word consists only of regular
local letters, and is classified by whether its length is at least two or is
one.  These alternatives are exhaustive.  They are disjoint because the class
is chosen by the first true predicate in a fixed ordered list.  Unique
historical payment of every primitive follows from the all-colour saturation
theorem, so classification cannot introduce a second charge.
\end{proof}

\subsection{Exact response projector and ledger}

Let $s=(s_0,\ldots,s_3)$ denote the four source variables on the fixed smaller
polydisc.  Let $\mathcal G$ be the finite-dimensional space of normalized
quadratic background germs and let
\[
 \Pm:\mathcal G\longrightarrow\mathbb R                    \tag{1.3}
\]
be the bounded transverse $F^2$ extractor used by compact V13--V15.  If
$\mathfrak L$ is the normalized logarithm of the exact pushed density, define
the required marginal response by
\[
 \cR_{\rm exact}:=
 \Pm\,\partial_{s_0}\partial_{s_1}\partial_{s_2}
              \partial_{s_3}\mathfrak L\big|_{s=0}.          \tag{1.4}
\]
Fewer physical marks are handled by deleting the unused derivatives; four are
kept because the imported normal-convergence theorem controls them.

Write $\mathfrak L_{\rm G}$ for the calibrated Gaussian-principal logarithm,
including its Haar-compensated determinant and declared finite jets, and put
\[
 \cR_{\rm main}:=
 \Pm\,\partial_{s_0}\partial_{s_1}\partial_{s_2}
              \partial_{s_3}\mathfrak L_{\rm G}\big|_{s=0}. \tag{1.5}
\]

\begin{theorem}[Exact non-double-counted remainder ledger]
\label{thm:exact-ledger}
On every finite-volume all-colour domain satisfying the imported V8 gate,
\[
\boxed{
 \cR_{\rm exact}=\cR_{\rm main}
 +E_{\rm vert}+E_{\rm polymer}+E_{\rm bad}
 +E_{\rm reset}+E_{\rm multiscale}.}                         \tag{1.6}
\]
The five error terms are the images under (1.4) of the five disjoint classes
of Lemma~\ref{lem:word-partition}, with
$E_{\rm multiscale}$ equal to the boundary class together with the endpoint of
the propagated stable history.  No determinant, local Taylor defect,
topology-changing exit, or reset complement occurs twice.
\end{theorem}

\begin{proof}
The exact triangular pushforward is the Gaussian calibrated term plus the
normally convergent sum of connected owner words.  Normal convergence on the
fixed source polydisc permits the four derivatives in (1.4) to pass through
that sum, and boundedness of $\Pm$ permits its application term by term.
Lemma~\ref{lem:word-partition} partitions the resulting series into five
disjoint subseries, proving (1.6).

The two-jet determinant is outside the owner algebra by the V6--V8 ownership
rule and therefore occurs only in $\cR_{\rm main}$.  A bad or sector-transition
word has one first historical bad birth and is placed wholly in $E_{\rm bad}$,
so none of its regular companions is recharged to $E_{\rm polymer}$.  A reset
word is similarly placed wholly in $E_{\rm reset}$.  Finally, only a genuine
global boundary tag and the terminal stable-history endpoint enter
$E_{\rm multiscale}$.  This proves the nonduplication statement.
\end{proof}

\begin{firewall}[What V1 does not prove]
Theorem~\ref{thm:exact-ledger} is an exact identity, not yet a smallness
estimate.  In particular, the reset and multiscale classes still require a
power-preserving propagation theorem.  No one-loop sign, continuum measure,
or Yang--Mills mass gap is claimed.
\end{firewall}

\begin{gate}[P1]
Gate P1 is closed: (1.6) is exhaustive and each contribution has one owner.
The next acquisition is a common $u^{6/5}$ bound for the local, polymer, and
bad classes in the already established analytic norm.
\end{gate}

\section{V2: local analytic, polymer, and bad-field response bounds}

\subsection{An explicit common primitive majorant}

Put
\[
 u=\beta^{-1},\qquad p=\frac65.                              \tag{2.1}
\]
The following notation only repackages the proved V6--V10 quantities:
\begin{align}
 a_3&=M_\ell(c_0/t_0)^3,
 &q_3(u)&=e^{a_3u^p}-1,                                      \tag{2.2}\\
 \epsilon_6(u)&=C_2q_3(u)
 +C_2 2^{d_*/2}e^{-c_0^2u^{-1/5}/32}
 +C_{\rm ext}u^{-d_*/2}e^{-c_0^2u^{-1/5}/2},                \tag{2.3}\\
 q_{\rm ker}(u)&=\frac{2\epsilon_6(u)}{1-\epsilon_6(u)},
 &q_{\rm sat}(u)&=
 \left(C_{\rm collar}u^{-d_*/2}e^{-c_{\rm col}/u}\right)^{1/M_{\rm sat}},
                                                                    \tag{2.4}\\
 q_{\rm comp}(u)&=e^{\tau_*C_{\rm comp}u^p}-1,              \tag{2.5}\\
 w(u)&=(K_{\chi_*}+1)\{q_{\rm ker}(u)+q_{\rm sat}(u)\}
       +q_{\rm comp}(u).                                     \tag{2.6}
\end{align}
Here $c_{\rm col}=\bar\epsilon_0^2/(6M_\kappa)$.  Formula (2.6)
is the $w_{10}$ strength, with every fresh regular, bad, and compensation
primitive paid once.

For $A,c,d,\beta_0>0$, define
\[
 \beta_\sharp(A,c,d;\beta_0)
 =\max\left\{\beta_0,\left(\frac{A}{cd}\right)^{1/d}\right\},
 \quad
 T(A,c,d;\beta_0)=\beta_\sharp^A e^{-c\beta_\sharp^d}.     \tag{2.7}
\]

\begin{lemma}[Explicit exponential-to-power conversion]
\label{lem:exp-power}
For $\beta\ge\beta_0$,
\[
 \beta^A e^{-c\beta^d}\le T(A,c,d;\beta_0).                \tag{2.8}
\]
Consequently each stretched-exponential term in (2.3)--(2.4) is
bounded by an explicit constant times $u^p$ on
$0<u\le u_0=\beta_0^{-1}$.
\end{lemma}

\begin{proof}
The logarithmic derivative of $f(\beta)=\beta^Ae^{-c\beta^d}$ is
$A/\beta-cd\beta^{d-1}$.  Thus $f$ increases up to
$(A/(cd))^{1/d}$ and decreases afterwards, proving (2.8).
To compare $\beta^ae^{-c\beta^d}$ with
$u^p=\beta^{-p}$, apply (2.8) with $A=a+p$.
\end{proof}

\begin{proposition}[Uniform $u^{6/5}$ owner strength]
\label{prop:w-power}
Choose $\beta_0$ so that
\[
 \epsilon_6(\beta_0^{-1})\le\frac12,\qquad
 a_3\beta_0^{-p}\le1,\qquad
 \tau_*C_{\rm comp}\beta_0^{-p}\le1.                      \tag{2.9}
\]
Then
\[
 \boxed{w(u)\le C_w(\beta_0)u^{6/5}}
 \qquad(0<u\le\beta_0^{-1}),                               \tag{2.10}
\]
where one explicit legal constant is
\begin{align}
 C_w={}&(K_{\chi_*}+1)[4C_6+C_{\rm sat}]
       +e\tau_*C_{\rm comp},                               \tag{2.11}\\
 C_6={}&e a_3C_2
 +C_2 2^{d_*/2}T(p,c_0^2/32,1/5;\beta_0)\\
 &+C_{\rm ext}T(d_*/2+p,c_0^2/2,1/5;\beta_0),              \tag{2.12}\\
 C_{\rm sat}={}&C_{\rm collar}^{1/M_{\rm sat}}
 T\left(\frac{d_*}{2M_{\rm sat}}+p,
         \frac{c_{\rm col}}{M_{\rm sat}},1;\beta_0\right).
                                                                    \tag{2.13}
\end{align}
\end{proposition}

\begin{proof}
For $0\le t\le1$, $e^t-1\le et$, so
$q_3\le ea_3u^p$ and
$q_{\rm comp}\le e\tau_*C_{\rm comp}u^p$.
Lemma~\ref{lem:exp-power} gives
$\epsilon_6\le C_6u^p$.  Since $\epsilon_6\le1/2$,
$q_{\rm ker}\le4C_6u^p$.  The same lemma, after the
$M_{\rm sat}$th root, gives
$q_{\rm sat}\le C_{\rm sat}u^p$.  Substitute in (2.6).
\end{proof}

\subsection{Local derivatives in the established norm}

Let $\|\cdot\|_{R,\boldsymbol\rho}^{\rm HK}$ be the compact
right-invariant Taylor--Wiener norm of compact V3, augmented by the source
radii, and let $0<r<R$.  Write $V_{\rm hi}$ for the regular local log-density
after exact extraction of its constant, linear, and quadratic field jets.

\begin{theorem}[High-grade local response]
\label{thm:local-response}
There is a fixed $C_{\rm loc}$ such that
\[
 \|V_{\rm hi}\|_{R,\boldsymbol\rho}^{\rm HK}
 \le C_{\rm loc}u^p,                                       \tag{2.14}
\]
and
\[
 \|DV_{\rm hi}\|_{r,\boldsymbol\rho}^{\rm HK}
 \le\frac{C_{\rm loc}}{R-r}u^p,\qquad
 \|D^2V_{\rm hi}\|_{r,\boldsymbol\rho}^{\rm HK}
 \le\frac{2C_{\rm loc}}{(R-r)^2}u^p.                      \tag{2.15}
\]
Restricting source radii from $\rho_i$ to $\rho_i'<\rho_i$ and
differentiating in a set $I$ of sources adds only
$\prod_{i\in I}(\rho_i-\rho_i')^{-1}$.
\end{theorem}

\begin{proof}
The V6 Morse identity writes the regular density as its exact two-jet times
$e^{\mathcal R_3}-1$; every other regular completion primitive is in
(2.5).  Proposition~\ref{prop:w-power}, with fixed owner embeddings, gives
(2.14).  The one-radius lemma for this same norm gives
$\|D^kF\|_r\le k!(R-r)^{-k}\|F\|_R$ for $k=1,2$.
Multivariable Cauchy gives the source statement.  No cutoff indicator is
differentiated.
\end{proof}

\subsection{Marked connected-polymer and bad-field bounds}

Let $M_{\rm rsp}$ contain one normalized marginal response mark, $\Pm$,
and the fixed owner/source embeddings.  Put
\[
 x(u)=A_\mu c_{\rm rsp}w(u),\qquad
 A_\mu=e^{1+\mu}(\Delta_*+1).                              \tag{2.16}
\]

\begin{theorem}[Source-marked remainder split]
\label{thm:marked-split}
If $x(u)<1$, the V1 subseries satisfy
\begin{align}
 |E_{\rm vert}|+|E_{\rm bad}|
 &\le C_{\rm vb}w(u),                                      \tag{2.17}\\
 |E_{\rm polymer}|
 &\le e^\mu M_{\rm rsp}c_{\rm rsp}w(u)
       \frac{2x(u)-x(u)^2}{(1-x(u))^2}.                     \tag{2.18}
\end{align}
If $x(u)\le1/2$, then
\[
 |E_{\rm vert}|+|E_{\rm bad}|+|E_{\rm polymer}|
 \le C_{\rm fresh}u^{6/5},                                \tag{2.19}
\]
where
\[
 C_{\rm fresh}=C_{\rm vb}C_w
 +8e^\mu M_{\rm rsp}c_{\rm rsp}^2A_\mu C_w^2u_0^p.       \tag{2.20}
\]
The constants are uniform in volume, colour stage, and regulator whenever the
imported local dictionary is uniform.
\end{theorem}

\begin{proof}
Singleton regular and bad words are bounded by their primitive variation
strength times the fixed response embeddings.  Historical saturation pays a
mixed word containing a bad primitive only to the bad class, proving (2.17).

For an all-regular connected word with $n\ge2$ primitives, root at the
response mark.  The first primitive costs
$e^\mu M_{\rm rsp}c_{\rm rsp}w$, each later primitive costs at most $x$,
and there are at most $n$ mark positions.  Hence
\[
 e^\mu M_{\rm rsp}c_{\rm rsp}w
 \sum_{n\ge2}n x^{n-1}
 =e^\mu M_{\rm rsp}c_{\rm rsp}w
   \frac{2x-x^2}{(1-x)^2},
\]
which proves (2.18).  For $x\le1/2$ the rational factor is at most $4x$.
Use (2.10) and $u^{2p}\le u_0^pu^p$ to obtain (2.19)--(2.20).
Normal convergence justifies the source derivatives.
\end{proof}

\begin{gate}[P2 and the local half of P3]
Theorem~\ref{thm:local-response} closes P2 in the existing strong/weak
analytic norms.  Theorem~\ref{thm:marked-split} closes the polymer and
response-weighted bad-field parts of P3.  The remaining task is propagation
of the reset/stable history without losing $p=6/5$.
\end{gate}

\section{V3: upstream reset propagation and multiscale ownership}

\subsection{Why the reset problem is a stable-history problem}

The V14 ledger of the main notes separates a same-scale Gaussian shell, a
local nonlinear correction, reference-reset memory, and a genuine global
boundary response.  A moment-preserving reset is not automatically harmless
at later scales.  However, after the constant, marginal two-jet, measured
redundant, anisotropy, and retained sector rows have been extracted, its
remaining local part lies in the stable owner space on which the archived
aligned map is contractive.  This is the correct upstream location of the
reset estimate.

Let $X_{\rm st}$ be that packet-polymer Banach space and let $s_j$ be the norm
of all propagated stable history entering scale $j$.  Fix constants
\[
 0\le q<1,\quad B_{\rm reset},L_h,M_c<\infty,              \tag{3.1}
\]
where $q$ is the cofinal stable contraction, $B_{\rm reset}$ is the norm of
stable projection followed by recharting, $L_h$ is the moving-reference
Lipschitz constant, and
\[
 |c_j|\le M_c,qquad
 \beta_{j+1}=\beta_j+c_j,qquad u_j=\beta_j^{-1}.           \tag{3.2}
\]
All are uniform finite-dictionary constants on the imported calibrated cone.

\begin{lemma}[Fresh reset forcing has the same power as the owner bank]
\label{lem:reset-forcing}
On the domain of Proposition~\ref{prop:w-power}, the stable forcing born
between scales $j$ and $j+1$ obeys
\[
 \|d_{{\rm st},j}\|_{X_{\rm st}}
 \le A_{\rm forc}u_j^p,                                    \tag{3.3}
\]
where one admissible constant is
\[
 A_{\rm forc}
 =B_{\rm reset}C_w+L_hM_cu_0^{,2-p}.                       \tag{3.4}
\]
\end{lemma}

\begin{proof}
Exact finite-jet extraction partitions the output into the calibrated jet and
its stable complement.  The local difference between exact and reference
outputs has owner norm at most $w(u_j)$; applying the bounded stable projection
and recharting costs $B_{\rm reset}$.  Proposition~\ref{prop:w-power} bounds
this part by $B_{\rm reset}C_wu_j^p$.

Changing the moving reference from $u_j$ to $u_{j+1}$ contributes at most
$L_h|u_{j+1}-u_j|$.  From (3.2),
\[
 |u_{j+1}-u_j|
 =\frac{|c_j|u_j^2}{|1+c_ju_j|}
 \le 2M_cu_j^2                                              \tag{3.5}
\]
when $M_cu_j\le1/2$.  Absorb the harmless factor $2$ into the declared
$L_h$, and use $u_j^2\le u_0^{2-p}u_j^p$.  The measured redundant projection
is exactly zero, retained labels are not reset, and bad histories were already
included in $w$; hence no omitted local term remains.
\end{proof}

\subsection{A power-preserving geometric convolution}

For $q>0$ put
\[
 \Lambda_q=\left(\frac{1+q}{2q}\right)^{1/p};
\]
for $q=0$ choose any $\Lambda_q>1$.  In either case
\[
 \theta_q:=q\Lambda_q^p<1.                                 \tag{3.6}
\]

\begin{theorem}[Reset-memory invariant cone]
\label{thm:reset-cone}
Assume
\[
 0<u_j\le u_*,\qquad
 M_cu_*\le\frac12,qquad
 1+M_cu_*\le\Lambda_q,                                    \tag{3.7}
\]
and the archived nonautonomous stable recursion
\[
 s_{j+1}\le q s_j+A_{\rm forc}u_j^p.                       \tag{3.8}
\]
Define
\[
 H_{\rm st}
 =\frac{A_{\rm forc}\Lambda_q^p}{1-q\Lambda_q^p}.         \tag{3.9}
\]
If $s_0\le H_{\rm st}u_0^p$, then for every $j\ge0$,
\[
 \boxed{s_j\le H_{\rm st}u_j^p.}                          \tag{3.10}
\]
Thus propagation does not degrade the exponent $p=6/5$.
\end{theorem}

\begin{proof}
Equation (3.2) gives
\[
 \frac{u_j}{u_{j+1}}=1+c_ju_j\le1+M_cu_*\le\Lambda_q,     \tag{3.11}
\]
while positivity of the couplings is part of the calibrated domain.
Suppose inductively that $s_j\le H_{\rm st}u_j^p$.  Then
\[
\begin{split}
 s_{j+1}
 &\le(qH_{\rm st}+A_{\rm forc})u_j^p\\
 &\le(qH_{\rm st}+A_{\rm forc})\Lambda_q^pu_{j+1}^p
 =H_{\rm st}u_{j+1}^p,
\end{split}                                                  \tag{3.12}
\]
where the last identity is (3.9).  The initial hypothesis starts the
induction.
\end{proof}

\begin{remark}[Noncircularity]
The hypothesis (3.8) is not the desired marginal estimate.  It is the archived
strict contraction theorem applied to the stable projection, with the fresh
forcing supplied independently by Lemma~\ref{lem:reset-forcing}.  The
marginal extractor never enters the proof of (3.10).
\end{remark}

\subsection{The genuine boundary class on the declared sector}

\begin{proposition}[No unowned boundary response on the fixed periodic sector]
\label{prop:boundary-zero}
For a periodic finite lattice, with flux, toron, and integer charge retained
as fixed labels and theta equal to zero, the boundary class in (1.2) has zero
local marginal response.  Every chart or admissibility exit is already in
$E_{\rm bad}$, and every surviving multiscale term is the stable-history
endpoint controlled by (3.10).
\end{proposition}

\begin{proof}
A periodic lattice has no spatial boundary face, so the sum of a discrete
divergence is zero.  The sectorwise pushforward theorem transports the flux,
toron, and integer charge labels instead of recomputing or summing them.
The V11 transition-ownership theorem assigns any path leaving an admissible
sector to its first historical collar-bad birth, already counted in
$E_{\rm bad}$.  The V8 saturation theorem assigns every chart-overlap failure
to the same unique bad history.  Hence the only endpoint not evaluated at its
birth scale is the stable reset history, which is (3.10).
\end{proof}

Let $L_{\rm st}$ be the uniform marked inverse-coupling response norm on
$X_{\rm st}$ supplied by the V13 marked expansion.

\begin{corollary}[Reset plus multiscale response]
\label{cor:reset-response}
Under Theorem~\ref{thm:reset-cone} and
Proposition~\ref{prop:boundary-zero},
\[
 |E_{\rm reset}|+|E_{\rm multiscale}|
 \le L_{\rm st}H_{\rm st}u_j^p.                           \tag{3.13}
\]
Together with (2.19), the complete non-Gaussian inverse-coupling error obeys
\[
 \boxed{|e_j|\le C_{\rm NP}u_j^{6/5}},\qquad
 C_{\rm NP}=C_{\rm fresh}+L_{\rm st}H_{\rm st}.          \tag{3.14}
\]
\end{corollary}

\begin{proof}
Integrate the bounded marked derivative along the segment from zero history
to the stable history of norm $s_j$, use (3.10), and add the fresh estimate
(2.19).  The owner partition in Theorem~\ref{thm:exact-ledger} makes the sum
exhaustive and disjoint.
\end{proof}

\begin{gate}[P3 and the inverse-coupling part of P4]
The reset convolution (3.10), the boundary classification, and (3.14) close
P3 on the fixed periodic retained sector and prove the critical order
separation in inverse-coupling variables: the main shell is order one while
the exhaustive nonperturbative error is $O(u^{6/5})$.  Iteration V4 converts
this into the exact physical-response margin and proves scale summability.
\end{gate}

\section{V4: physical margin protection and summation over scales}

\subsection{Exact reciprocal conversion}

Let $b_j$ be the main programme's certified Gaussian-principal
inverse-coupling response at scale $j$.  Define $e_j$ by
\[
 c_j=b_j+e_j,\qquad
 u_{j+1}=\frac{u_j}{1+u_jc_j}.                              \tag{4.1}
\]
Thus $c_j$ is the exhaustive exact inverse-coupling increment and, by
Corollary~\ref{cor:reset-response},
\[
 |e_j|\le C_{\rm NP}u_j^p,\qquad p=\frac65.              \tag{4.2}
\]
Assume the already established bounded-increment card
\[
 |c_j|\le M_c.                                               \tag{4.3}
\]
No numerical value of $b_j$ is used below.

Define the exact and main physical responses by
\[
 \cR_j^{\rm exact}=u_{j+1}-u_j,\qquad
 \cR_j^{\rm main}=-b_ju_j^2.                               \tag{4.4}
\]

\begin{lemma}[Exact physical remainder identity]
\label{lem:physical-identity}
Whenever $1+u_jc_j\ne0$,
\[
 \cR_j^{\rm exact}-\cR_j^{\rm main}
 =-e_ju_j^2+\frac{c_j^2u_j^3}{1+u_jc_j}.                   \tag{4.5}
\]
If $M_cu_j\le1/2$, then
\[
 \boxed{
 |\cR_j^{\rm exact}-\cR_j^{\rm main}|
 \le C_{\rm NP}u_j^{16/5}+2M_c^2u_j^3.}                   \tag{4.6}
\]
\end{lemma}

\begin{proof}
For every scalar $c$ with $1+uc\ne0$,
\[
 \frac{u}{1+uc}=u-cu^2+\frac{c^2u^3}{1+uc}.
\]
Insert $c_j=b_j+e_j$ and subtract $u_j-b_ju_j^2$, proving
(4.5).  Under $M_cu_j\le1/2$ the denominator in the last term has
absolute value at least $1/2$.  Apply (4.2)--(4.3).
\end{proof}

\subsection{The symbolic-margin theorem}

The only datum requested from the main line is a lower margin
\[
 |b_j|\ge m_{\rm main}>0.                                  \tag{4.7}
\]
All other constants below are outputs of the remainder programme.

Let $u_{\rm owner}$ be the minimum of the thresholds required for
Proposition~\ref{prop:w-power}, $x(u)\le1/2$, and the imported calibrated
cone.  Put
\[
\boxed{
 u_*=
 \min\left\{
 u_{\rm owner},
 \frac1{2M_c},
 \frac{\Lambda_q-1}{M_c},
 \left(\frac{m_{\rm main}}{4C_{\rm NP}}\right)^{5/6},
 \frac{m_{\rm main}}{8M_c^2}
 \right\}.}                                                \tag{4.8}
\]
If $C_{\rm NP}=0$, the fourth entry is interpreted as $+\infty$.
The third entry is omitted when $q=0$ and $\Lambda_q$ was chosen without a
restriction.  Shrinking $u_{\rm owner}$ once also ensures the source radii
and every strict analytic inclusion used above.

\begin{theorem}[Regulator-uniform nonperturbative margin protection]
\label{thm:margin-protection}
On every fixed-sector periodic admissible weak-coupling trajectory with
$0<u_j<u_*$, (4.7) implies
\[
\boxed{
 |\cR_j^{\rm exact}-\cR_j^{\rm main}|
 \le
 \left[
 \frac{C_{\rm NP}}{m_{\rm main}}u_j^{6/5}
 +\frac{2M_c^2}{m_{\rm main}}u_j
 \right]|\cR_j^{\rm main}|
 \le\frac12|\cR_j^{\rm main}|.}                          \tag{4.9}
\]
Consequently
\[
 \operatorname{sgn}\cR_j^{\rm exact}
 =\operatorname{sgn}\cR_j^{\rm main}.                    \tag{4.10}
\]
The estimate is uniform in finite volume, colour stage, regulator, and scale
on the declared trajectory.
\end{theorem}

\begin{proof}
By (4.7),
$|\cR_j^{\rm main}|=|b_j|u_j^2\ge m_{\rm main}u_j^2$.
Divide (4.6) by this lower bound.  The fourth and fifth entries of (4.8)
make the two displayed summands in (4.9) at most $1/4$ each.
Thus the distance from the exact real response to the nonzero main response
is less than its magnitude, which proves (4.10).  Uniformity follows from the
uniform constants in (2.10), (2.19), (3.10), and (4.3).
\end{proof}

\begin{corollary}[One-step nonlinear RG remainder]
\label{cor:one-step}
The exact physical marginal recursion has the form
\[
 u_{j+1}=u_j-b_ju_j^2+r_j,                                   \tag{4.11}
\]
with
\[
 \boxed{|r_j|\le C_{\rm NP}u_j^{16/5}+2M_c^2u_j^3
 =O(u_j^3).}                                                 \tag{4.12}
\]
In particular $r_j=o(u_j^2)$, so Gate P4 is positive rather than a failure
stop.
\end{corollary}

\begin{proof}
Equation (4.11) is the definition (4.4); (4.12) is (4.6).
Since $16/5>3>2$, both terms are strictly higher order than the leading
response.
\end{proof}

\subsection{Scale-summable accumulation}

For accumulation in the ultraviolet orientation, suppose in addition that
\[
 b_j\ge m_{\rm main}>0.                                    \tag{4.13}
\]
The threshold (4.8) makes
$|e_j|\le m_{\rm main}/4$, hence
\[
 c_j\ge\frac34m_{\rm main}\ge\frac12m_{\rm main}.       \tag{4.14}
\]

\begin{lemma}[Explicit power sum along the protected trajectory]
\label{lem:power-sum}
If $\beta_{j+1}\ge\beta_j+a$ with $a>0$, then for every $\alpha>1$,
\[
 \sum_{j=0}^\infty u_j^\alpha
 \le \beta_0^{-\alpha}
 +\frac{\beta_0^{1-\alpha}}{a(\alpha-1)}.                \tag{4.15}
\]
\end{lemma}

\begin{proof}
Induction gives $\beta_j\ge\beta_0+aj$.  The function
$(\beta_0+at)^{-\alpha}$ is decreasing, so its sum is bounded by its
value at zero plus its integral on $[0,\infty)$.  Evaluation of the
integral gives (4.15).
\end{proof}

\begin{theorem}[Summable nonlinear response error]
\label{thm:summable}
Under (4.13) and $u_0<u_*$, the protected domain is forward invariant and
\[
 \sum_{j=0}^\infty
 |\cR_j^{\rm exact}-\cR_j^{\rm main}|<\infty.            \tag{4.16}
\]
More precisely,
\begin{equation}
\begin{split}
 E_{\rm NP}^{\rm total}
 :={}&\sum_{j\ge0}
 |\cR_j^{\rm exact}-\cR_j^{\rm main}|\\
 \le{}&
 C_{\rm NP}\left[
 \beta_0^{-16/5}
 +\frac{10}{11m_{\rm main}}\beta_0^{-11/5}\right]\\
 &+2M_c^2\left[
 \beta_0^{-3}+\frac1{m_{\rm main}}\beta_0^{-2}\right].
\end{split}
\tag{4.17}
\end{equation}
\end{theorem}

\begin{proof}
Equation (4.14) gives
$\beta_{j+1}\ge\beta_j+m_{\rm main}/2$, so $u_j$ decreases and can never
leave $0<u<u_*$.  Apply Lemma~\ref{lem:power-sum} to (4.6), first with
$a=m_{\rm main}/2$, $\alpha=16/5$, and then with $\alpha=3$.
The two integral coefficients are
$2/[m_{\rm main}(11/5)]=10/(11m_{\rm main})$ and
$2/[m_{\rm main}(2)]=1/m_{\rm main}$, giving (4.17).
\end{proof}

\begin{gate}[P4 and P5]
Corollary~\ref{cor:one-step} closes the critical go/no-go gate P4 with a
strict power separation.  Theorem~\ref{thm:summable} closes P5.  Theorem
\ref{thm:margin-protection} already has the requested success-stop form, with
an explicit $u_*$ and only the symbolic margin $m_{\rm main}$ supplied by the
main programme.  Iteration V5 performs the mandatory companion sweep and
checks whether any new note invalidates or strengthens an imported premise
before the stop is declared.
\end{gate}

\section{V5: mandatory companion sweep and success stop}

\subsection{Files inspected}

The fifth-iteration sweep was performed after V4.  The exact snapshot was
\[
\begin{array}{l|r|r|l}
\text{file}&\text{lines}&\text{bytes}&\text{SHA--256 prefix}\\ \hline
\texttt{Renewal\_Yang\_Mills\_Mass\_Gap\_Research\_Notes\_V77.tex}
&26825&913923&\texttt{6CA0089783961EDA}\\
\texttt{Renewal\_SM\_Einstein\_Continuum\_Research\_Notes\_V26.tex}
&8127&288229&\texttt{4C4BF2D7A80239EB}\\
\texttt{Renewal\_NS\_Stability\_Research\_Notes\_V26.tex}
&5319&278283&\texttt{82C887F8C2C69E4F}.
\end{array}                                                   \tag{5.1}
\]
The full hashes are retained in the filesystem audit used to form this note.

The Yang--Mills stream advanced from V74 to V77 during V1--V4.  V75 proves a
typed cyclic trace identity, V76 finds that a factorwise differentiated
majorant is worse than the operative V69 bound, and V77 introduces
energy-dressed trace words but obtains only floating finite-grid diagnostics
for the relevant scale.  Each of V75--V77 explicitly states that it does not
provide a higher-loop or nonperturbative estimate.  These results do not alter
(3.14): the present bound is an absolute owner majorant and uses no
factorwise trace cancellation.  The typed-cycle lesson reinforces the decision
not to split the main Gaussian response before its packet-first recombination.

The SM--Einstein stream advanced from V21 to V26.  Its V26 normal
commuting-square limit shows how a vanishing residual can pass through
interacting conditional expectations.  No Yang--Mills conditional-expectation
square with the required owner-norm estimate is supplied there, so importing
that result would add an unproved hypothesis.  It is unnecessary because the
strict stable recursion already proves (3.10).  The Navier--Stokes file is
unchanged and contains no compact determinant, sector boundary, or YM reset
estimate.

\begin{record}[V5 audit decision]
No new companion result invalidates an imported premise or improves the
nonperturbative exponent.  The correct direction remains the one taken in
V3--V4: complete-word ownership, stable reset propagation, and reciprocal
conversion.  No redirection is required.
\end{record}

\subsection{The five-gate ledger}

\begin{center}
\begin{tabular}{c|p{0.62\textwidth}|c}
gate&result&status\\ \hline
P1&Exact disjoint owner identity (1.6).&closed\\
P2&Uniform high-grade local response and two field derivatives,
      (2.14)--(2.15).&closed\\
P3&Source-marked connected polymer, bad-field, reset, and
      multiscale bounds, (2.19) and (3.13).&closed\\
P4&Exact one-step remainder is $O(u^3)=o(u^2)$, (4.12).&go\\
P5&The accumulated physical-response error is summable,
      (4.16)--(4.17).&success
\end{tabular}
\end{center}

\subsection{Final theorem in the requested interface form}

\begin{theorem}[Success-stop theorem]
\label{thm:success-stop}
Let the finite-volume exact Yang--Mills block trajectory be periodic, remain in
the calibrated weak-coupling cone, retain its flux, toron, and integer-charge
labels, and satisfy the already-proved all-colour owner gate and cofinal stable
contraction.  Let the main programme provide only the certified number
$m_{\rm main}>0$ such that
$|b_j|\ge m_{\rm main}$ on the same trajectory.  For $0<u_j<u_*$, with
$u_*$ given explicitly by (4.8),
\[
\boxed{
 \left|\cR_j^{\rm exact}-\cR_j^{\rm certified\ main}\right|
 \le E_{\rm NP}(u_j)
 :=C_{\rm NP}u_j^{16/5}+2M_c^2u_j^3
 \le\frac12\left|\cR_j^{\rm certified\ main}\right|.}  \tag{5.3}
\]
If $u_*\le1$, then also
\[
 E_{\rm NP}(u)
 \le
 \frac{C_{\rm NP}+2M_c^2}{m_{\rm main}},
 u\,|\cR^{\rm certified\ main}(u)|,                       \tag{5.4}
\]
so the requested relative form holds with $\delta=1$.  In the positive
ultraviolet orientation the errors are scale summable with the explicit total
bound (4.17).  Therefore the exact and certified-main responses have the same
sign at every scale in the domain.
\end{theorem}

\begin{proof}
The exact owner identity is Theorem~\ref{thm:exact-ledger}.  Its fresh
local, polymer, and bad subseries satisfy (2.19); reset and multiscale history
satisfy (3.13); hence their exhaustive inverse-coupling sum satisfies (3.14).
Lemma~\ref{lem:physical-identity} converts that estimate without a truncated
formal expansion.  The threshold (4.8) and
Theorem~\ref{thm:margin-protection} give (5.3) and the sign conclusion.
For $u\le1$, $u^{6/5}\le u$, so division of (4.6) by the main lower bound
$m_{\rm main}u^2$ gives (5.4).  The scale sum is
Theorem~\ref{thm:summable}.
\end{proof}

\begin{firewall}[Exact claim boundary at the stop]
The theorem protects the sign of a certified packet-first marginal response
on the declared weak-coupling trajectory.  It does not supply the numerical
margin $m_{\rm main}$, certify the still-open Brillouin quadrature in V77,
construct an infinite-volume or continuum Yang--Mills measure, prove
Osterwalder--Schrader reconstruction, establish exponential clustering, or
prove a Yang--Mills mass gap.  Those are outside this parallel programme.
\end{firewall}

\begin{gate}[SUCCESS STOP]
The success criterion is reached by (5.3)--(5.4): the exhaustive nonlinear
remainder is regulator-uniform, source differentiable, strictly higher order,
uniformly less than one half of the symbolic main margin for explicit $u_*$,
and summable over the ultraviolet trajectory.  In accordance with the stated
stopping rule, the programme stops here; no constant optimization is pursued.
\end{gate}

\clearpage
\section{Programme II: Yang--Mills sector tightness and topological
transport}
\label{sec:sector-programme}

\subsection{Context, target, and nonduplication boundary}

Programme I is complete and is not reopened below.  The new question is
whether the finite-cutoff Yang--Mills laws can satisfy, simultaneously in
cutoff and finite volume,
\[
 \sup_j\mu_j\{\Psi(q_j)\ge R\}\le C e^{-cR},\qquad c>0,       \tag{6.1}
\]
together with coherent transport of $q_j$ under refinement.
\emph{Volume-uniform} is understood in the locality-compatible sense: the
class of finite regulators is closed under finite disjoint unions, on which
the Wilson action, product Haar measure, and exact block integrations
factorize.

The programme memorandum fixes the gates and stopping rules; it is not
evidence for a mathematical assertion.  The following results are imported
without reproof.
\begin{enumerate}[label=\textup{(S\arabic*)},leftmargin=4em]
\item Main YM V11 separates local $\operatorname{Tr}F\widetilde F$, integer
      global charge, and retained flux/toron data.  Its extended block map
      copies the global labels and its regular fibres are sector diagonal.
\item Archived YM V27 proves Haar-preserving erasure and an exponential area
      bound for detail-exact centre sheets, expressly excluding winding
      sheets.
\item Archived YM V28 proves a volume-uniform multi-plaquette bound and a
      rooted noncentral bad-cluster tail for the bare Wilson law.  These are
      local/rooted estimates, not a tail for the total number of defects.
\item Archived YM V45 proves that torons and boundary holonomies are compact
      external slow variables, not residual tails.
\item Companion SM V74--V78 proves the abstract energy--entropy and gapped
      refinement implications, but supplies no Yang--Mills constants.
\end{enumerate}
The task is therefore to decide whether those hypotheses can hold for a
global, volume-uniform Yang--Mills sector variable.

\subsection{S1: exact sector ledger for the retained block map}

Fix a four-dimensional periodic torus and $G=SU(3)$.  Write
\[
 q_\Lambda=(z_\Lambda,\xi_\Lambda,Q_\Lambda),                \tag{6.2}
\]
where
\[
 z_\Lambda\in H^2(\mathbb T^4;\mathbb Z_3)
                 \simeq(\mathbb Z_3)^6,
\]
and where the primitive noncontractible holonomies lie in the compact
simultaneous-conjugacy quotient
\[
 {\cal X}_{\rm hol}:=G^4/G.                                  \tag{6.3}
\]
Its closed commuting locus is the toron space.  The coordinate
$\xi_\Lambda$ may also retain the finite boundary-holonomy data of the
chosen tree gauge.  On the gauge-invariant admissible open set
$\cA_\Lambda$, the geometric charge is the continuous integer-valued map
$Q_\Lambda:\cA_\Lambda\to\mathbb Z$ of main YM V11.  Use a cemetery value
$\partial$ off $\cA_\Lambda$ when a total measurable label is needed.

\begin{lemma}[Only the integer coordinate can escape]
\label{lem:only-integer-escapes}
The product
\[
 {\cal K}_{\rm top}
 :=H^2(\mathbb T^4;\mathbb Z_3)\times{\cal X}_{\rm hol}
\]
is compact.  Every family of probability laws on its coordinates is tight.
Thus a proper sector complexity need only be unbounded in $|Q_\Lambda|$;
adding a bounded continuous function of $(z_\Lambda,\xi_\Lambda)$ changes
only the constants in (6.1).
\end{lemma}

\begin{proof}
The cohomology group is finite.  The group $G^4$ is compact, and its quotient
by the continuous action of compact $G$ is compact.  Their product is a
common compact containment set.  Every continuous real function on it is
bounded.
\end{proof}

\begin{proposition}[Exact preservation and the only mixing mechanisms]
\label{prop:sector-ledger}
For the actual extended block map of main YM V11,
\[
 \widehat{\cal B}_\Lambda(U)
 =
 \bigl({\cal B}_\Lambda U,z_\Lambda(U),
       \xi_\Lambda(U),Q_\Lambda(U)\bigr),                    \tag{6.4}
\]
the ledger is:
\begin{enumerate}[label=\textup{(\roman*)},leftmargin=3.5em]
\item cubical subdivision preserves the centre-flux class through its
      induced isomorphism on $H^2$;
\item tree gauge retains primitive noncontractible chords, and coarse
      ordered-path multiplication preserves their holonomies up to
      base-point conjugacy;
\item $Q_\Lambda$ is constant on each connected regular fibre and is copied,
      not recomputed, by (6.4);
\item labels mix only if a coordinate is forgotten, recomputed from the
      coarse local field, or crosses $\cA_\Lambda^c$.  The last event is
      already assigned to the saturated-bad owner.  If fine sectors merge,
      their probability-weighted conditional mixture, not each individual
      conditional, is the coarse conditional law.
\end{enumerate}
\end{proposition}

\begin{proof}
Subdivision is a homotopy equivalence and induces a cohomology isomorphism.
Fine-path multiplication telescopes around the coarse cycle, while changing
base point only conjugates it.  A continuous integer-valued map on a
connected fibre is constant.  Formula (6.4) then gives exact retention.
The listed operations exhaust the failures of that argument, and ordinary
disintegration proves the weighted-mixture statement.
\end{proof}

\begin{gate}[S1]
Finite centre flux and compact holonomy/toron coordinates are already tight
and exactly retained.  The only noncompact global label is integer charge.
Winding-sheet erasure and toron concentration are therefore not substitutes
for an integer-charge estimate.
\end{gate}

\subsection{S2: energy information and its exact ceiling}

There is no positive Wilson-action penalty for moving through the commuting
toron space.

\begin{proposition}[Commuting torons have zero Wilson action]
\label{prop:toron-zero-action}
Let $H_1,\ldots,H_4\in SU(3)$ commute.  On a periodic lattice with side
lengths $L_\mu$, there is a spatially constant link field whose primitive
$\mu$-cycle holonomy is $H_\mu$ and whose every plaquette is the identity.
Thus its Wilson action is zero.
\end{proposition}

\begin{proof}
The commuting tuple lies in a common maximal torus.  The power map
$T\mapsto T^{L_\mu}$ is surjective on that torus, so choose commuting
$T_\mu$ with $T_\mu^{L_\mu}=H_\mu$ and set $U_{x,\mu}=T_\mu$.
Every plaquette is
$T_\mu T_\nu T_\mu^{-1}T_\nu^{-1}=I$, whereas the primitive cycle product is
$T_\mu^{L_\mu}=H_\mu$.
\end{proof}

This zero is harmless because toron space is compact.  For centre-labelled
contractible sheets, archived V27 already gives the linear area cost
\[
 \bigl[\beta\sigma_0(\delta)-\kappa_\Psi\bigr]|\Sigma|.
\]
For noncentral local defects, archived V28 gives the stronger multi-event
probability bound.  Neither result controls a winding cohomology class or
global integer charge.

Let a finite connected component $\Lambda_0$ have unnormalised sector weights
\[
 W(q)=\int_{\{Q=q\}}e^{-S}\,dH.                              \tag{6.5}
\]
The identity neighbourhood gives $W(0)>0$.  If a nonzero admissible sector is
nonempty, it is open in the admissible set.  Product Haar measure gives it
positive measure and the finite Wilson density is strictly positive, so
$W(q_*)>0$ for some $q_*\ne0$.

\begin{lemma}[Additivity caps a global charge penalty at linear growth]
\label{lem:energy-linear-ceiling}
On the disjoint union of $N$ copies of $\Lambda_0$, let action and charge be
additive and Haar measure be the product.  Normalize the nonnegative Wilson
action to have global minimum zero in the identity, charge-zero sector.  If a
charge-$q_*$ configuration has finite action $E_*$, then for $0\le k\le N$
the minimum action at total charge $kq_*$ obeys
\[
 0\le E_N(kq_*)\le kE_*.                                    \tag{6.6}
\]
Thus no local additive action can force a superlinear penalty in the number
of separated charge carriers.
\end{lemma}

\begin{proof}
Place the charge-$q_*$ configuration on $k$ components and the zero-action
identity configuration on the remaining $N-k$.  Additivity gives action
$kE_*$.  Taking the infimum in the charged sector can only lower it, while
nonnegativity gives the other inequality.
\end{proof}

\begin{gate}[S2]
The local energy bounds are retained, but the global integer-charge penalty
has an exact linear ceiling under separation.  Granting a linear lower bound
would be the strongest locality-compatible rate; S3 must compare it with
placement entropy.
\end{gate}

\subsection{S3: exact sector-placement entropy}

\begin{theorem}[Binomial sector multiplicity]
\label{thm:binomial-sector-entropy}
The unnormalised weight on $N$ disjoint copies of configurations having
exactly $k$ components in sector $q_*$ and all others in sector zero is
\[
 \binom Nk W(q_*)^kW(0)^{N-k}.                              \tag{6.7}
\]
Its placement entropy satisfies
\[
 S_N(k):=\log\binom Nk\ge k\log\frac Nk
 \qquad(1\le k\le N).                                      \tag{6.8}
\]
If
\[
 \Delta_*=-\log\frac{W(q_*)}{W(0)}
\]
is the complete one-component sector free-energy cost, already including
internal entropy, then
\[
 k\Delta_*-S_N(k)
 \le k\left(\Delta_*-\log\frac Nk\right)
 \longrightarrow-\infty                                   \tag{6.9}
\]
for every fixed $k\ge1$ as $N\to\infty$.
\end{theorem}

\begin{proof}
Factorization gives weight $W(q_*)^kW(0)^{N-k}$ for each choice of the $k$
charged components.  The choices are disjoint and number $\binom Nk$.
Moreover
\[
 \binom Nk
 =\prod_{i=0}^{k-1}\frac{N-i}{k-i}
 \ge\left(\frac Nk\right)^k,
\]
because $(N-i)/(k-i)\ge N/k$ when $N\ge k$.  This proves (6.8);
substitution gives (6.9).
\end{proof}

\begin{gate}[S3]
The entropy count is exact.  Even after replacing bare action by the full
one-component sector free energy, placement entropy is unbounded with volume
at fixed complexity.  Hence the net sector-weight exponent and comparison
floor required by companion V76 cannot both be volume independent for a
global extensive charge.
\end{gate}

\subsection{S4: the volume-uniform global tail is impossible}

If charge is not defined almost surely, include
$\mathbf 1_{\cA_\Lambda^c}$ in the nonnegative sector complexity.  On
disjoint copies, any positive one-component inadmissibility probability
already gives the Bernoulli divergence in the first part of the next proof.
Thus it remains only to treat an almost-sure integer charge.

\begin{theorem}[Tensorization no-go for global sector tightness]
\label{thm:global-sector-no-go}
Let $X$ be the integer charge under the normalized Yang--Mills law on
$\Lambda_0$, and suppose $X$ is not almost surely zero.  On $N$ disjoint
copies let $X_1,\ldots,X_N$ be the independent charges and define
\[
 \Psi_N=\sum_{i=1}^N|X_i|,
 \qquad
 Q_N=\sum_{i=1}^NX_i.                                      \tag{6.10}
\]
For every fixed $R<\infty$,
\[
 \mathbb P\{\Psi_N\ge R\}\longrightarrow1.                 \tag{6.11}
\]
The signed total-charge laws are not tight either:
\[
 \mathbb P\{|Q_N|\ge R\}\longrightarrow1.                  \tag{6.12}
\]
Consequently no $C<\infty$ and $c>0$, independent of volume, can bound
either tail by $Ce^{-cR}$ for every $N$ and $R$.
\end{theorem}

\begin{proof}
Put $p=\mathbb P\{|X|\ge1\}>0$.  The indicators
$B_i={\bf1}_{\{|X_i|\ge1\}}$ are independent Bernoulli variables of
parameter $p$, and $\Psi_N\ge\sum_iB_i$.  The weak law gives
$N^{-1}\sum_iB_i\to p$ in probability, proving (6.11).

The charge $X$ is bounded on the fixed finite lattice, hence has finite
variance.  If $m=\mathbb EX\ne0$, then $Q_N/N\to m$ in probability and
(6.12) follows.  If $m=0$, nontriviality gives
$\sigma^2=\operatorname{Var}X>0$.  The central limit theorem gives
$Q_N/(\sigma\sqrt N)\Rightarrow Z$, with $Z$ standard normal.  For every
$\eta>0$ and all sufficiently large $N$,
\[
 \mathbb P\{|Q_N|\le R\}
 \le\mathbb P\{|Q_N|/(\sigma\sqrt N)\le\eta\}.
\]
The portmanteau bound for the closed interval $[-\eta,\eta]$, followed by
$\eta\downarrow0$, proves that this probability tends to zero.  In the
remaining zero-variance case $X$ is constant; if it is not zero, the
nonzero-mean argument applies.

Finally choose $R_0$ with $Ce^{-cR_0}<1/2$.  Equations (6.11) and (6.12)
make the corresponding probability exceed $1/2$ for large $N$, a
contradiction.
\end{proof}

\begin{corollary}[The energy--entropy failure criterion is reached]
\label{cor:failure-stop}
For any nontrivial global charge sector,
\[
 \boxed{E(R)-S_{\rm entropy}(R)\not\longrightarrow+\infty}
                                                                  \tag{6.13}
\]
uniformly in volume.  More sharply, (6.9) tends to $-\infty$ along fixed
$R=k$ and increasing volume.  The proposed energy--entropy mechanism cannot
prove (6.1).
\end{corollary}

\begin{proof}
Lemma~\ref{lem:energy-linear-ceiling} gives the additive energy ceiling and
Theorem~\ref{thm:binomial-sector-entropy} gives the placement entropy.  Their
difference is (6.9), while the probability-level contradiction is
Theorem~\ref{thm:global-sector-no-go}.
\end{proof}

\begin{firewall}[Connected tori versus the mechanism no-go]
\label{firewall:connected-volume}
Theorem~\ref{thm:global-sector-no-go} is an exact obstruction to the stated
volume-uniform theorem on a locality-compatible regulator class, because
local Yang--Mills measures tensorize on disjoint unions.  It does not by
itself prove anti-concentration of total charge along a sequence restricted
only to connected periodic tori.  Such a theorem would require positive
topological susceptibility or a uniform lower insertion/mixing estimate,
neither of which occurs in the current notes.  Restricting to fixed physical
volume can also remove placement divergence.  Those are different targets;
they do not rescue the volume-uniform global energy--entropy mechanism.
\end{firewall}

\subsection{Refinement transport retained at the failure boundary}

Although S4 terminates the tightness route, its independent refinement
statement can be recorded exactly.  Let $P_j:X_{j+1}\to X_j$ be the field
block map, $r_j:{\cal S}_{j+1}\to{\cal S}_j$ the label map, and let
$\mathcal G_{j+1}$ be the set on which
\[
 r_jq_{j+1}=q_jP_j.
\]
For probability laws $\mu_j$, put
\[
 \epsilon_j=\mu_{j+1}({\cal G}_{j+1}^c),\qquad
 d_j=\|(P_j)_*\mu_{j+1}-\mu_j\|_{\rm TV}.                  \tag{6.14}
\]

\begin{proposition}[Refinement defect is data-processing stable]
\label{prop:refinement-defect}
The label laws obey
\[
 \left\|(r_j)_*(q_{j+1})_*\mu_{j+1}
       -(q_j)_*\mu_j\right\|_{\rm TV}
 \le\epsilon_j+d_j.                                        \tag{6.15}
\]
For the copied labels in (6.4), $\epsilon_j=0$ on the extended state.
If $\sum_j(\epsilon_j+d_j)<\infty$, successive label laws are projectively
Cauchy after pushforward to every fixed level.
\end{proposition}

\begin{proof}
Couple $r_jq_{j+1}(x)$ and $q_jP_j(x)$ using the same
$x\sim\mu_{j+1}$.  They differ only off $\mathcal G_{j+1}$, so their
total-variation distance is at most $\epsilon_j$.  Applying $q_j$ cannot
increase total variation; hence the distance from
$(q_jP_j)_*\mu_{j+1}$ to $(q_j)_*\mu_j$ is at most $d_j$.  The triangle
inequality proves (6.15), and summation and telescoping give the last claim.
\end{proof}

\begin{remark}[Coherence does not prevent escape]
Point masses $\delta_j$ on $\mathbb Z$ may be perfectly coherent under a
changing-label convention while escaping every finite sector set.
Proposition~\ref{prop:refinement-defect} preserves labels faithfully; it
does not make an extensive global label tight.
\end{remark}

\subsection{Companion checkpoint and terminal decision}

The checkpoint immediately before this decision inspected
\[
\begin{array}{l|r|l}
\text{file}&\text{bytes}&\text{SHA--256 prefix}\\ \hline
\texttt{Renewal\_Yang\_Mills\_Mass\_Gap\_Research\_Notes\_V82.tex}
 &958424&\texttt{D4E3B41C26CF2E1C}\\
\texttt{Renewal\_SM\_Einstein\_Continuum\_Research\_Notes\_V39.tex}
 &455984&\texttt{9B0F17930C29371F}\\
\texttt{Renewal\_NS\_Stability\_Research\_Notes\_V26.tex}
 &278283&\texttt{82C887F8C2C69E4F}
\end{array}                                                   \tag{6.16}
\]
Main YM V81--V82 add a signed-diagonal frequency collapse and a rational
path-cell certificate for the Brillouin programme.  They add no sector-energy
or topology estimate; the inherited V35 warning that local plaquette control
is not global sector tightness remains operative.  SM V38 proves a
projectively compatible fixed point for compact affine holonomy, reinforcing
the compact-coordinate treatment in S1, while SM V39 concerns an
OS-recovery isometry and vacuum multiplicity.  Neither changes the global
integer-charge entropy.  The NS file is unchanged.  The archived centre
erasure, chessboard, and toron-patch results therefore remain the strongest
relevant Yang--Mills inputs; none defeats tensorization.

\begin{center}
\begin{tabular}{c|p{0.62\textwidth}|c}
gate&result&status\\ \hline
S1&Finite flux and compact holonomy labels separated from integer charge;
exact extended-map preservation.&closed\\
S2&Local centre/noncentral estimates retained; toron action is flat and
global charge energy has an additive linear ceiling.&closed\\
S3&Exact binomial placement entropy, (6.7)--(6.9).&closed\\
S4&Every nontrivial extensive charge law violates a volume-uniform
exponential global tail under tensorization.&failure\\
S5&Label coherence obeys (6.15), but cannot repair escape.&structural only\\
S6&Joint continuum handoff.&not entered
\end{tabular}
\end{center}

\begin{gate}[FAILURE STOP: energy--entropy tightness route fails]
\label{gate:sector-failure-stop}
The programme's stated failure criterion is reached:
\[
 \boxed{\text{energy--entropy tightness route fails}.}       \tag{6.17}
\]
The obstruction is not a poor constant.  Global integer charge and absolute
topological activity are extensive, additive variables; separated carriers
have a linear energy/free-energy cost and unbounded placement entropy.  A
cutoff- and volume-uniform tail for their total value is therefore the wrong
infinite-volume target unless all nonzero sectors have zero probability.

The appropriate return to the main YM programme is to retain compact
flux/toron labels exactly, use fixed-window or charge-density estimates for
local topology, and assemble the continuum through phase decomposition or a
quasi-local observable algebra.  A fixed-physical-volume charge tail is also
a legitimate separate programme.  In accordance with the stop rule, no OS
reconstruction, clustering, mass-gap argument, or S6 continuum construction
is attempted here.
\end{gate}

\end{document}
